\subsection{Problèmes recontrés}

La première difficulté que nous pouvons citer porte sur l'émetteur-récepteur d'ultrasons.
Malheureusement, l'algorithme utilisé ne parvenait pas à sortir d'une boucle infinie ce qui rendait le Robot incapable de suivre les autres programmes.
Nous avons tenté de paralléliser ce dispositif mais sans succès, c'est pourquoi nous n'avons pas pu présenter cette option lors de la soutenance.\\

Pendant une grande partie du projet, nous avions le sentiment d'avoir une longueur de retard sur les autres groupes à propos de la partie électronique. 
Ce retard était dû aux nombreuses discussions que nous avions eu en amont sur la conception du Robot.
Nous avons discuté du nombre de capteurs, de la position des moteurs et d'autres dispositifs pendant plusieurs séances avant de commencer à percer et à souder.
Au final, cela ne nous a pas desservi; effectivement, après le montage, tous le câblage rentrait parfaitement dans le châssis et la plupart des composants ont fonctionné du premier coup.
De plus, la plupart des erreurs rencontrées venaient des codes utilisés.

Cependant, lors de nos premiers esssais, les moteurs ne répondaient pas. C'était une simple erreur de câblage sans gravité (aucun matériel n'a été dégradé), 
mais avec la fatigue, nous avons mal vérifié l'ensemble des connectiques et le problème ne fût résolu que lors de la séance d'essais suivante.

Avec ce retard accumulé, nous avons manqué de temps pour les derniers réglages, c'est pourquoi un dernier bug subsistait lors de la présentation du Robot.\\

